% =============================================================================
% Section 5 -- Formal APIP Schema
% Contains Table 2 (tab:schema), typeset as a full-width table* so the field
% descriptions stay readable.
% =============================================================================
\section{Formal APIP Schema}
\label{sec:schema}

We formalize the APIP as a declarative schema to enable tooling integration and
audit trail generation. The schema uses a notation analogous to JSON Schema;
YAML, JSON, or DSL implementations are straightforward.

\begin{table*}[t]
  \centering
  \caption{Formal APIP schema fields. New fields relative to a basic PIP
           schema are \texttt{kirkpatrick\_levels}, the
           \texttt{exogenous\_disruption} failure mode,
           \texttt{mesh\_context}, and the attribution-provenance fields
           \texttt{attribution\_confidence}, \texttt{mapping\_table\_hash},
           and \texttt{environment\_provenance}.}
  \label{tab:schema}
  \footnotesize
  \begin{tabular}{@{}L{4.2cm}L{1.8cm}L{10.4cm}@{}}
    \toprule
    \textbf{Field} & \textbf{Type} & \textbf{Description} \\
    \midrule
    \texttt{apip\_id} & string &
      Unique identifier for this APIP instance. \\
    \texttt{agent\_id} & string &
      Identifier of the agent subject to this APIP. \\
    \texttt{created\_at} & datetime &
      Timestamp when the APIP was initiated. \\
    \texttt{triggered\_by} & object &
      The metric, threshold value, and observed value that triggered
      the APIP. \\
    \texttt{behavioral\_contract\_ref} & string &
      Reference to the behavioral contract document active at deployment
      time. \\
    \texttt{failure\_mode} & enum &
      One of: \texttt{behavioral\_drift} $|$ \texttt{input\_brittleness} $|$
      \texttt{alignment\_creep} $|$ \texttt{tool\_misuse} $|$
      \texttt{exogenous\_disruption} $|$ \texttt{abstain}, plus the
      out-of-scope extension value \texttt{adversarial\_subversion}
      (Section~\ref{sec:taxonomy}). \\
    \texttt{cause\_attribution\_report} & string &
      Reference to the structured cause attribution document. \\
    \texttt{attribution\_confidence} & number &
      Confidence score emitted by the attribution funnel
      (Section~\ref{subsubsec:funnel}); \texttt{abstain} is recorded when
      below the frozen threshold. \\
    \texttt{mapping\_table\_hash} & string &
      Hash of the frozen trace-label-to-taxonomy mapping table active for
      this APIP, making post-hoc edits to attribution rules detectable. \\
    \texttt{environment\_provenance} & object &
      Version pins for the agent's environment at trigger time (model
      fingerprints, knowledge-base/corpus versions, external dependency
      commits), copied into the record rather than referenced. \\
    \texttt{kirkpatrick\_levels} & object &
      Baseline and target values for Levels 1--4 evaluation at each
      check-in. \\
    \texttt{intervention\_tier} & enum &
      One of: \texttt{prompt\_level} $|$ \texttt{retrieval\_level} $|$
      \texttt{model\_level}. \\
    \texttt{interventions} & array &
      List of specific interventions applied, each with type, description,
      and timestamp. \\
    \texttt{window\_start} & datetime &
      Start of the remediation window. \\
    \texttt{window\_end} & datetime &
      Scheduled end of the remediation window. \\
    \texttt{checkins} & array &
      Scheduled check-in timestamps and recorded metric snapshots with
      Kirkpatrick level assessments. \\
    \texttt{exit\_outcome} & enum &
      One of: \texttt{resolved} $|$ \texttt{regressed} $|$ \texttt{timeout}
      $|$ \texttt{recalibrated} $|$ \texttt{retired}
      (Section~\ref{subsec:exit}). \\
    \texttt{exit\_evaluated\_at} & datetime &
      Timestamp of the formal exit evaluation. \\
    \texttt{owner} & string &
      Identity of the human operator or team responsible for this APIP. \\
    \texttt{mesh\_context} & object &
      If applicable: \texttt{mesh\_id}, \texttt{peer\_agents},
      \texttt{context\_weights} at trigger time, exogenous flag. \\
    \texttt{notes} & string &
      Free-text field for qualitative observations throughout the
      lifecycle. \\
    \bottomrule
  \end{tabular}
\end{table*}

The schema is intentionally minimal at the field level but requires structured
sub-schemas for \texttt{interventions} and \texttt{checkins} to preserve audit
trail integrity. The \texttt{behavioral\_contract\_ref} field is
architecturally critical: it enforces the discipline of writing a behavioral
contract at deployment time, without which the APIP trigger and exit evaluation
phases cannot be operationalized. The \texttt{kirkpatrick\_levels} field
captures the multi-level evaluation data at each check-in, enabling post-hoc
analysis of whether interventions produced superficial or durable improvement.
The \texttt{mapping\_table\_hash} and \texttt{environment\_provenance} fields
extend the visibility mechanisms of Kolt et al.~\cite{kolt2024visibility} from
identity linkage to \emph{decision-rule} linkage: an auditor can verify not
only which agent and operator acted, but that the attribution rules and
environment versions cited by the record are the ones that were actually in
force---which is what converts a pre-registration claim from an assertion into
a checkable property.
The \texttt{mesh\_context} field, described further in
Section~\ref{subsec:mesh}, supports cause attribution in multi-agent
deployments.
